\section{Plant Phenotyping}

\paragraph{}
Plant breeding is defined as the process of inducing desirable inheritable traits in plants through controlled human action  \cite{ortonIntroduction2020}. This process is important in solving the problem of adequate global food supply. The rapid growth of the human population has created a strain on global food supply. It is estimated that crop production must increase by 70\% by 2050 \cite{carvajal-yepesGlobalSurveillanceSystem2019} to meet this demand. Additionally, abiotic and biotic stressors limit cultivar biomass and yield \cite{grimeEvidenceExistenceThree1997}. Therefore, plant breeders need effective techniques to produce germsplam for cultivars that maximise yield and biomass, in various adverse environmental conditions. In the last few decades, image processing techniques in plant phenotyping have been shown to be an effective tool for plant breeding \cite{anandBreedingTechniquesDispense2023}.

\subsection{What Is Plant Phenotyping?}
Plant phenotyping is the evaulation of phenotypes in plants \cite{fioraniFutureScenariosPlant2013}. Phenotypes are primitive or complex traits of an organism that can be observed or quantified and evaulated numerically \cite{johannsenGenotypeConceptionHeredity1911a}.  The literature also indicates that phenotype can alternatively be seen as combination of all morphological, physiological, chemical, and behavioral aspects that represent the individual plant \cite{guoPhenotypingPlants2014, mohammedReviewPlantPhenotyping, liReviewImagingTechniques2014}. 

\subsection{Plant Phenotyping As Method For Improving Plant Traits}
Quantative data produced by plant phenotyping is most useful in analyzing plant-environment interactions in plant breeding \cite{fioraniFutureScenariosPlant2013} for development of germsplasm with improved plant traits \cite{liReviewImagingTechniques2014}. Phenotypical data is usefuf as a proxy for assesing plant performance in relation to its environment. For example, stressors or disturbances of plants are induced by environmental conditions that organisms are placed in, and these stressors induce symptoms which can be detected and evaulated using phenotyping. This indicates phenotypical traits of a plant have complex dynamic feedback loop with their environment \cite{walterPlantPhenotypingBean2015}, and therefore phenotypical traits serve as proxies in determining if certain genotypes have favourable traits in particular environmental conditions.

\subsection{Image-Obtainable Phenotypical Traits As Indicators Of Plant Health, Growth, And Stress}
Phenotypical traits that can be obtained using image processing techniques can be divided into morphological, physiological, and temporal traits. Additionally the literature indicates, morphological and physiological measurements are classifiable as holistic-based or component-based \cite{daschoudhuryLeveragingImageAnalysis2019}. Holostic traits consider the plant as a whole, whereas component-based features look at particular organs like the fruit, flower, leaves, stem etc. of the organism. 

\paragraph{Leaf phenotypinc traits} \hfill \\
Leaf-derived traits are one of the most commonly used type of traits, because leaves rapidly respond to environmental stimuli, and robustly indicate various factors which determine plant health, growth and stress. \cite{zhangHighthroughputPhenotypingPlant2023}. Leaf phenotypic traits have been used to predict plant health, growth and stress. For example, the morphological trait of leaf area predicts plant productivity, because reduced leaf area results in decreased chlorophyll content, and consequently less photosynthesis. Therefore leaf area is a good proxy measure for evaluating plant growth, and has been used to monitor and estimate crop yield, and health \cite{zhangLeafAreaIndex2021}. 

%Everthing below is part of draft
 
 \subparagraph{Morphological traits}  About plant sturcture, which affects its performance and gives indication regarding its health
 \begin{itemize}
 	\item Leaf width
 	\item Number of leaves
 	\item Leaf inclination angle
 	\item Leaf area
 	\item Root system architecture (e.g root diameter) 
 	\item Plant shoots
 	\item Fruit
 	\item Flowers
 \end{itemize}
 
 \subparagraph{Physiological traits}  About how plant operates in repsonse to its growing conditions
 \begin{itemize}
 	\item Stomotal conductance
 	\item Chlorophyll content of leaf
 	\item Plant water status
 	\item Plant photosynthesis (metablism)
 	\item Nutrient content
 \end{itemize}
 
 \subparagraph{Phenological traits} These are phenological traits (when certain biological events occur that depend on climates)
 \begin{itemize}
 	\item When fruit or flowers are yielded
 	\item Yield prediction
 \end{itemize}
 
\paragraph{The use of plant phenotyping and detecting plant stress}
Plant biomass is affected by stresses and disturbances \cite{grimeEvidenceExistenceThree1997}.
A stressor is anything which stymies the growth of plant biomass. Examples of stressors include things like water scarcity, inadequate amounts of light, and unfavourable temperatures.  A disturbance is anything which causes some level of destruction to the plant. Examples of disturbances include pathogens, fires, frost, and pests

\textbf{Biotic Stress + Disturbance}.
\begin{itemize}
	\item Pathogenic organisms
	\item Pests
	\item Parasitc plants
\end{itemize}

\textbf{Abiotic Stress + Disturbance}
\begin{itemize}
	\item Drought
	\item Flood
	\item Salinity
	\item Heat Stress
	\item Raditation Stress, light stress
\end{itemize}